\documentclass[11pt,twocolumn]{article}

% ---------- Geometry ----------
\usepackage[margin=1in]{geometry}

% ---------- Math ----------
\usepackage{amsmath,amssymb,amsthm,mathtools}

% ---------- Typography ----------
\usepackage{microtype}
\usepackage[T1]{fontenc}

% ---------- Tables & Figures ----------
\usepackage{booktabs}
\usepackage{graphicx}
\usepackage{float}
\usepackage{caption}

% ---------- Algorithms ----------
\usepackage[ruled,vlined,linesnumbered]{algorithm2e}

% ---------- Colors & Links ----------
\usepackage[dvipsnames]{xcolor}
\usepackage[colorlinks=true,linkcolor=RoyalBlue,citecolor=OliveGreen,urlcolor=BrickRed]{hyperref}
\usepackage{cleveref}

% ---------- Lists ----------
\usepackage{enumitem}

% ---------- Bibliography ----------
\usepackage{natbib}
\bibliographystyle{plainnat}

% ---------- Theorem environments ----------
\newtheorem{theorem}{Theorem}[section]
\newtheorem{lemma}[theorem]{Lemma}
\newtheorem{proposition}[theorem]{Proposition}
\newtheorem{corollary}[theorem]{Corollary}
\theoremstyle{definition}
\newtheorem{definition}[theorem]{Definition}
\newtheorem{remark}[theorem]{Remark}

% ---------- Custom commands ----------
\newcommand{\E}{\mathbb{E}}
\newcommand{\Var}{\operatorname{Var}}
\newcommand{\R}{\mathbb{R}}
\newcommand{\N}{\mathbb{N}}
\newcommand{\Aspace}{\mathcal{A}}
\newcommand{\Astar}{\mathcal{A}^{\ast}}
\newcommand{\Pspace}{\bar{\mathcal{P}}}
\newcommand{\pref}{p_{\text{ref}}}
\newcommand{\pLLM}{p_{\text{LLM}}}
\newcommand{\Qdirect}{Q_{\text{direct}}}
\newcommand{\bx}{\mathbf{x}}
\newcommand{\bn}{\mathbf{n}}

% ================================================================
\title{\textbf{EveryStep: Path Tracing the Agentic Workflow}\\[6pt]
  \large Applying Monte Carlo Light Transport Theory \\
  to LLM Agent Reliability Optimization}

\author{%
  (Authors)\\
  \textit{Working Paper --- March 2026}
}

\date{}

% ================================================================
\begin{document}
\maketitle

% ================================================================
\begin{abstract}
Large Language Models are often dismissed as unreliable because each call is stochastic.
We argue that this critique, while valid for single invocations, fundamentally misunderstands the problem:
\emph{path tracing in computer graphics faces the exact same challenge and has solved it.}
We propose a formal mapping between Monte Carlo light transport algorithms and LLM agentic workflows,
grounded in three insights:
(1)~LLM unreliability is a variance problem amenable to 60+ years of graphics variance reduction;
(2)~the \emph{Point Light Theorem}: tasks whose correct answers occupy a zero-measure subset of the output space
    require Next Event Estimation (NEE), mapping directly to RAG and tool use;
(3)~a \emph{sampling efficiency} criterion $\eta_M = -\ln(1-p_M)/C_M$ determines when cheaper models
    with more samples outperform expensive models.
We derive the \emph{Agent Quality Equation} as an analogue of the rendering equation,
formalize six variance reduction techniques
(NEE, MIS, Russian Roulette, Importance Sampling, MLT, path-space MIS),
prove two theorems on the necessity of NEE for ``point light'' tasks,
and design experiments to validate the framework on Art Bible generation.
\end{abstract}

% ================================================================
\section{Introduction}\label{sec:intro}

\subsection{The ``Gacha'' Critique}

A prevalent argument holds that because each LLM output is sampled from a distribution,
the result is inherently unreliable.
If a single call succeeds with probability $p$,
trusting it means accepting failure rate $1-p$.
This is correct---for a single sample.

\subsection{Path Tracing: The Same Problem, Solved}

Path tracing computes global illumination by shooting random rays into a scene.
A single sample per pixel is noise.
Yet modern renderers produce photorealistic images because
\begin{equation}\label{eq:failure}
  P(\text{failure after } N \text{ samples}) = (1-p)^N.
\end{equation}
For $p=0.5$, $N=10$: $(0.5)^{10} \approx 10^{-3}$, i.e.\ 99.9\% success.
But path tracing goes far beyond naive repetition.
Over six decades, the rendering community developed
importance sampling, MIS~\citep{veach1995optimally},
NEE, Russian Roulette, and MLT~\citep{veach1997metropolis}---each
provably improving convergence under specific conditions.

\subsection{Core Thesis and Research Questions}

We propose that Monte Carlo light transport is \emph{directly transferable}
to LLM agentic workflows:
\begin{enumerate}[label=\textbf{RQ\arabic*},leftmargin=2em]
  \item Can path tracing algorithms systematically improve industry-level agent performance?
  \item Which rendering concepts map onto which agent concepts, and where does the analogy break?
  \item Can variance reduction enable cheaper models to outperform expensive ones at fixed budget?
  \item Are these optimizations generalizable across modalities and task types?
\end{enumerate}

\subsection{Art Bible as Target Application}

An Art Bible is a comprehensive visual style guide for a creative project.
It is a long-chain stochastic process (world-building $\to$ color $\to$ material $\to$ character $\to$ \ldots),
each step depending on all predecessors.
Correctness is defined by a \emph{range} (not a point), analogous to diffuse surfaces.
Global consistency is required, analogous to energy conservation.

% ================================================================
\section{Mathematical Foundations}\label{sec:math}

\subsection{The Rendering Equation}

\citet{kajiya1986rendering} introduced the rendering equation:
\begin{equation}\label{eq:rendering}
  L_o(\bx, \omega_o) = L_e(\bx, \omega_o)
    + \int_\Omega f_r(\bx, \omega_i, \omega_o)\,
      L_i(\bx, \omega_i)\,
      (\omega_i \cdot \bn)\,d\omega_i
\end{equation}
where $L_o$ is outgoing radiance, $L_e$ is emission,
$f_r$ is the BRDF, $L_i$ is incoming radiance,
and $(\omega_i \cdot \bn)$ the cosine term over hemisphere $\Omega$.

\subsection{The Agent Quality Equation}

We propose the \emph{Agent Quality Equation} as a direct analogue:
\begin{equation}\label{eq:agent-quality}
  Q(s_k) = \Qdirect(s_k)
    + \int_{\Aspace} T(s_k, a)\,Q(s_{k+1}(a))\,p(a \mid s_k)\,da
\end{equation}
where:
\begin{itemize}[nosep]
  \item $Q(s_k)$: total quality at state $s_k$ \quad($\leftrightarrow L_o$)
  \item $\Qdirect(s_k)$: direct contribution \quad($\leftrightarrow L_e$)
  \item $T(s_k, a)$: transfer function \quad($\leftrightarrow f_r$)
  \item $Q(s_{k+1}(a))$: quality from continuation \quad($\leftrightarrow L_i$)
  \item $p(a \mid s_k)\,da$: LLM probability measure \quad($\leftrightarrow (\omega_i\!\cdot\!\bn)\,d\omega_i$)
  \item $\Aspace$: action space \quad($\leftrightarrow \Omega$)
\end{itemize}

The recursive structure is identical: a second-kind Fredholm integral equation
where total quality equals direct contribution plus the integral over all continuations,
weighted by transfer and probability.

\subsection{Path Integral Formulation}

Following \citet{veach1997robust}, we reformulate in path space.
A rendering path $\bar{x} = (x_0, \ldots, x_K)$ contributes:
\begin{equation}\label{eq:veach-path}
  I = \int_{\bar\Omega} f(\bar{x})\,d\mu(\bar{x})
\end{equation}

An \emph{agent execution path}
$\bar{p} = (s_0, a_0, s_1, a_1, \ldots, s_K)$ contributes:
\begin{equation}\label{eq:agent-path}
  \mathcal{Q} = \int_{\Pspace} F(\bar{p})\,d\nu(\bar{p})
\end{equation}
where
\begin{align}
  F(\bar{p}) &= Q_{\text{final}}(\bar{p}) \cdot \prod_{k=0}^{K-1} T(s_k, a_k) \label{eq:path-throughput} \\
  d\nu(\bar{p}) &= \prod_{k=0}^{K-1} p(a_k \mid s_k)\,da_k \label{eq:path-measure}
\end{align}

\subsection{Monte Carlo Estimation}\label{sec:mc-basics}

The basic MC estimator for $\mathcal{Q}$:
\begin{equation}\label{eq:mc-estimator}
  \hat{\mathcal{Q}}_N = \frac{1}{N}\sum_{i=1}^{N} \frac{F(\bar{p}_i)}{p(\bar{p}_i)}
\end{equation}
with $\bar{p}_i \overset{\text{iid}}{\sim} p$.
The variance is
\begin{equation}\label{eq:mc-variance}
  \Var[\hat{\mathcal{Q}}_N]
    = \frac{1}{N}\,\Var\!\left[\frac{F(\bar{p})}{p(\bar{p})}\right]
    = \frac{\sigma^2}{N}
\end{equation}

\begin{remark}
Variance converges as $O(1/N)$.
Halving the standard deviation requires $4\times$ samples.
All techniques in \cref{sec:algorithms} aim to reduce $\sigma^2$,
thereby achieving better convergence per sample.
\end{remark}

\subsection{Sampling Efficiency Criterion}\label{sec:efficiency}

Let model $M$ have per-sample cost $C_M$ and success probability $p_M$.
After $N$ samples:
\begin{equation}\label{eq:success-prob}
  P_{\text{success}}(N) = 1 - (1-p_M)^N
\end{equation}

Under budget $B$, affordable samples: $N = \lfloor B/C_M \rfloor$.
Small model $M_s\;(p_s, C_s)$ dominates large $M_l\;(p_l, C_l)$ when:
\begin{equation}\label{eq:dominance}
  (1-p_s)^{B/C_s} < (1-p_l)^{B/C_l}
\end{equation}
Taking logarithms and rearranging:

\begin{definition}[Sampling Efficiency]\label{def:efficiency}
\begin{equation}\label{eq:eta}
  \boxed{\eta_M = \frac{-\ln(1-p_M)}{C_M}}
\end{equation}
When $p_M \ll 1$, $\eta_M \approx p_M / C_M$.
The model with higher $\eta$ is preferred at any budget.
\end{definition}

This formalizes the finding of \citet{brown2024monkeys}: 5 small-model samples
can outperform 1 GPT-4o sample because $\eta_{\text{small}} > \eta_{\text{large}}$.

% ================================================================
\section{Core Algorithms}\label{sec:algorithms}

\subsection{Next Event Estimation (NEE)}\label{sec:nee}

\subsubsection{The Point Light Problem in Rendering}

A point light occupies zero solid angle:
\begin{equation}\label{eq:point-light-measure}
  \mu\bigl(\{\omega \in \Omega : \omega \text{ hits point light}\}\bigr) = 0
\end{equation}
Pure path tracing \emph{never} finds it.
NEE solves this by explicit light sampling:
\begin{equation}\label{eq:nee-rendering}
  \hat{L}_{\text{NEE}}(\bx) = \sum_{l \in \mathcal{L}}
    \frac{f_r(\bx, \omega_l, \omega_o)\,L_e(\bx_l)\,G(\bx, \bx_l)}{p_{\text{light}}(\bx_l)}
    \cdot V(\bx, \bx_l)
\end{equation}
where $G(\bx,\bx_l) = \frac{|\cos\theta_x||\cos\theta_l|}{\|\bx - \bx_l\|^2}$
is the geometry term and $V \in \{0,1\}$ the visibility (shadow ray).

\subsubsection{Agent NEE}

``Point lights'' in agentic workflows are tasks with
$\mu(\{a \in \Aspace : a \text{ is correct}\}) \approx 0$:
exact API parameters, strict JSON schemas, precise computations.

\emph{NEE-Agent} explicitly connects to reference sources:
\begin{equation}\label{eq:nee-agent}
  \hat{Q}_{\text{NEE}}(s_k)
    = \Qdirect(s_k)
    + \frac{T(s_k, a^*)\,Q^*(s_{k+1})}{\pref(a^*)}
      \cdot \mathbb{1}[\operatorname{valid}(a^*, s_k)]
\end{equation}
where $a^* \sim \pref$ (RAG, tool call, verified exemplar)
and $\mathbb{1}[\operatorname{valid}(\cdot)]$ acts as the shadow ray.

\subsection{Multiple Importance Sampling (MIS)}\label{sec:mis}

\subsubsection{Balance Heuristic}

Given $n$ strategies with $n_i$ samples from $p_i$
\citep{veach1995optimally}:
\begin{equation}\label{eq:mis}
  \hat{F}_{\text{MIS}} = \sum_{i=1}^{n} \frac{1}{n_i}
    \sum_{j=1}^{n_i} w_i(X_{i,j})\,\frac{f(X_{i,j})}{p_i(X_{i,j})}
\end{equation}
with the balance heuristic weight:
\begin{equation}\label{eq:balance}
  w_i(x) = \frac{n_i\,p_i(x)}{\sum_{k=1}^{n} n_k\,p_k(x)}
\end{equation}

\begin{theorem}[Veach's Variance Bound {\citep{veach1995optimally}}]\label{thm:veach-bound}
For the balance heuristic with $\sum_i n_i = N$:
\begin{equation}\label{eq:veach-bound}
  \Var[\hat{F}_{\mathrm{MIS}}]
    \leq \frac{1}{\min_i n_i}
    \left(\int |f(x)|\,dx\right)^{\!2}
\end{equation}
\end{theorem}

The \emph{power heuristic} ($\beta = 2$) often improves upon this:
\begin{equation}\label{eq:power-heuristic}
  w_i^{(\beta)}(x) = \frac{(n_i\,p_i(x))^\beta}{\sum_{k=1}^{n}(n_k\,p_k(x))^\beta}
\end{equation}

\subsubsection{Agent MIS}

Three fundamental strategies:
\begin{enumerate}[nosep]
  \item $p_1$: Free LLM generation (BSDF sampling)
  \item $p_2$: RAG-guided generation (light sampling / NEE)
  \item $p_3$: Tool/API calls (delta distribution)
\end{enumerate}
\begin{equation}\label{eq:agent-mis}
  \hat{Q}_{\text{MIS}}(s_k) = \sum_{i=1}^{3} \frac{1}{n_i}
    \sum_{j=1}^{n_i}
    \frac{n_i\,p_i(a_{i,j})}{\sum_{m=1}^{3} n_m\,p_m(a_{i,j})}
    \cdot \frac{F(a_{i,j})}{p_i(a_{i,j})}
\end{equation}

MIS automatically allocates credit to the best strategy per sample---handling
the entire task spectrum from ``diffuse'' to ``point light'' without manual tuning.

\subsection{Russian Roulette}\label{sec:rr}

At bounce $k$ with throughput $\beta_k$:
\begin{equation}\label{eq:rr}
  \hat{L} = \begin{cases}
    L_{\text{cont}} / q_k & \text{w.p.\ } q_k \\
    0 & \text{w.p.\ } 1-q_k
  \end{cases}
\end{equation}
where $q_k = \min(\|\beta_k\|, 1)$.

\begin{proposition}[Unbiasedness of Russian Roulette]\label{prop:rr-unbiased}
$\E[\hat{L}] = q_k \cdot \frac{L}{q_k} + (1-q_k) \cdot 0 = L$.
\end{proposition}

\subsubsection{Agent Russian Roulette}

A Validator Agent evaluates partial quality $Q_k \in [0,1]$:
\begin{equation}\label{eq:agent-rr}
  q_k = \min\!\left(\frac{Q_k}{Q_{\text{threshold}}},\,1\right)
\end{equation}
Low-quality paths terminate early; survivors are upweighted by $1/q_k$.
The variance increase is bounded:
\begin{equation}\label{eq:rr-var-bound}
  \Var_{\text{RR}} \leq \Var_{\text{no-RR}} \cdot \left(\frac{1}{q_k} - 1\right)
\end{equation}
which is small for $q_k \approx 1$ (good paths).

\subsection{Importance Sampling}\label{sec:is}

The zero-variance IS PDF:
\begin{equation}\label{eq:optimal-is}
  q^*(x) = \frac{|f(x)|\,p(x)}{\int |f(x')|\,p(x')\,dx'}
\end{equation}

\subsubsection{Agent IS via Constraint Buffer}

Extract constraints $\mathcal{C}_k = \{c_1, \ldots, c_m\}$ from step $k$,
then bias step $k{+}1$:
\begin{equation}\label{eq:guided}
  p_{\text{guided}}(a \mid s_{k+1})
    \propto \pLLM(a \mid s_{k+1}) \cdot \prod_{j=1}^{m} \phi(a, c_j)
\end{equation}
where $\phi(a, c_j) \geq 0$ is a compatibility score.
This concentrates sampling in high-quality regions,
approximating importance sampling.

\subsection{Metropolis Light Transport (MLT)}\label{sec:mlt}

Given path $\bar{p}$ with quality $F(\bar{p})$,
propose mutation $\bar{p}' \sim T(\bar{p}' \mid \bar{p})$.
Accept with \citet{metropolis1953equation,hastings1970monte} probability:
\begin{equation}\label{eq:mh}
  \alpha(\bar{p}' \mid \bar{p})
    = \min\!\left(1,\,
      \frac{F(\bar{p}')\,T(\bar{p} \mid \bar{p}')}{F(\bar{p})\,T(\bar{p}' \mid \bar{p})}
    \right)
\end{equation}
For symmetric proposals:
\begin{equation}\label{eq:mh-symmetric}
  \alpha = \min\!\left(1,\,\frac{F(\bar{p}')}{F(\bar{p})}\right)
\end{equation}

\subsubsection{Agent MLT: Mutation-Based Document Refinement}

\begin{algorithm}[H]
\SetAlgoLined
\KwIn{Initial draft $\bar{p}_0$, iterations $M$, quality function $F$}
\KwOut{Refined document $\bar{p}_{\text{best}}$}
$\bar{p}_{\text{best}} \leftarrow \bar{p}_0$\;
\For{$t = 1$ \KwTo $M$}{
  Select section index $j$ uniformly at random\;
  $\bar{p}' \leftarrow \text{Mutate}(\bar{p}_{\text{current}}, j)$ \tcp*{LLM revises section $j$}
  $\alpha \leftarrow \min(1,\; F(\bar{p}') / F(\bar{p}_{\text{current}}))$\;
  \With{probability $\alpha$}{
    $\bar{p}_{\text{current}} \leftarrow \bar{p}'$\;
    \If{$F(\bar{p}') > F(\bar{p}_{\text{best}})$}{
      $\bar{p}_{\text{best}} \leftarrow \bar{p}'$\;
    }
  }
}
\Return $\bar{p}_{\text{best}}$\;
\caption{Agent MLT for Document Refinement}
\label{alg:mlt}
\end{algorithm}

The Markov chain converges to a stationary distribution $\pi(\bar{p}) \propto F(\bar{p})$,
spending more time in high-quality document regions.

% ================================================================
\section{The Point Light Theorem}\label{sec:point-light}

This section contains the sharpest theoretical contribution.

\begin{definition}[Agent Point Light]\label{def:point-light}
A task exhibits the \emph{point light property} if the correct output set
$\Astar \subset \Aspace$ satisfies
\begin{equation}
  \mu(\Astar) = 0
\end{equation}
under the LLM's output measure $\mu$ induced by $p(\cdot \mid s)$.

More practically, a task is \emph{$\epsilon$-point-light} if
$\mu(\Astar) < \epsilon$ for $\epsilon \ll 1$.
\end{definition}

\begin{definition}[Task Reflectance Spectrum]\label{def:spectrum}
Tasks are classified along a continuous spectrum:

\begin{center}
\begin{tabular}{@{}lll@{}}
\toprule
BRDF Type & $\mu(\Astar)$ & Example \\
\midrule
Diffuse & $\approx \mu(\Aspace)$ & ``Write a poem'' \\
Glossy & $\ll \mu(\Aspace)$ & ``Write an Art Bible'' \\
Specular & single point & ``$\int_0^1 x^2\,dx = \;?$'' \\
Point light & $= 0$ & Exact API call \\
\bottomrule
\end{tabular}
\end{center}
\end{definition}

\begin{theorem}[Unreachability by Pure Sampling]\label{thm:unreachable}
For a task with the point light property,
the naive MC estimator $\hat{Q}_N$ satisfies:
\begin{equation}\label{eq:unreachable}
  P\!\left(\exists\, i \leq N : a_i \in \Astar\right) = 0
  \quad \forall\, N \in \N
\end{equation}
when $a_i \sim p(\cdot \mid s)$ and $\mu(\Astar) = 0$.
\end{theorem}

\begin{proof}
Each $a_i$ is drawn independently from $p(\cdot \mid s)$,
which induces measure $\mu$.
Since $\mu(\Astar) = 0$, we have $P(a_i \in \Astar) = 0$ for each $i$.
By countable subadditivity:
\begin{equation*}
  P\!\left(\bigcup_{i=1}^{N} \{a_i \in \Astar\}\right)
    \leq \sum_{i=1}^{N} P(a_i \in \Astar) = 0
    \qedhere
\end{equation*}
\end{proof}

\begin{theorem}[NEE Reachability]\label{thm:nee-reach}
If there exists a reference distribution $\pref$ with
$\pref(\Astar) > 0$, then the NEE-Agent estimator satisfies:
\begin{equation}\label{eq:nee-reach}
  \E\bigl[\hat{Q}_{\text{NEE}}\bigr] > 0
  \quad \text{and} \quad
  \Var\bigl[\hat{Q}_{\text{NEE}}\bigr] < \infty
\end{equation}
\end{theorem}

\begin{proof}
The NEE estimator draws $a^* \sim \pref$ and evaluates
$T(s, a^*)\,Q^*(s')/\pref(a^*)$.
Since $\pref(\Astar) > 0$, the event $\{a^* \in \Astar\}$ has positive probability.
For correct actions, $T(s, a^*) > 0$ and $Q^*(s') > 0$, so:
\begin{equation*}
  \E\bigl[\hat{Q}_{\text{NEE}}\bigr]
    \geq \pref(\Astar) \cdot \inf_{a^* \in \Astar}
      \frac{T(s, a^*)\,Q^*(s')}{\pref(a^*)} > 0
\end{equation*}
For finite variance, we require $T \cdot Q^* / \pref$ to be square-integrable
w.r.t.\ $\pref$.
Since $\pref$ has support on $\Astar$ and $T, Q^*$ are bounded,
$\E[(T \cdot Q^* / \pref)^2] < \infty$.
\end{proof}

\begin{corollary}[Necessity of NEE]\label{cor:nee-necessary}
For tasks with $\mu(\Astar) = 0$,
NEE-style direct sampling (RAG, tool calls, verified references)
is a \textbf{necessary condition} for non-zero expected quality.
No amount of random sampling suffices.
\end{corollary}

\subsection{The MIS Bridge}\label{sec:mis-bridge}

In rendering, MIS handles mixed scenes (diffuse + point lights) by combining
BSDF sampling with NEE.
The balance heuristic automatically zeros contributions from strategies
assigning zero probability.

\begin{proposition}[Agent MIS handles the full spectrum]\label{prop:mis-spectrum}
For a task with both ``diffuse'' components (many valid outputs)
and ``point light'' components (exact answers),
the Agent MIS estimator~\eqref{eq:agent-mis} with strategies
$p_1$ (free), $p_2$ (RAG), $p_3$ (tool)
satisfies:
\begin{equation}
  \E[\hat{Q}_{\text{MIS}}] = \mathcal{Q}
\end{equation}
(unbiased), with variance no greater than $O(1/\min_i n_i)$ times the
integral of $|F|$ squared (\cref{thm:veach-bound}).
\end{proposition}

\begin{proof}
Unbiasedness follows from the MIS weights summing to 1 for each sample point:
$\sum_i w_i(x) = 1$.
The variance bound is a direct application of \cref{thm:veach-bound}.
For point-light components, $p_1(a^*) = 0$ so $w_1(a^*) = 0$;
all weight falls on $p_3$ (tools), exactly as light sampling dominates
for point lights in rendering.
\end{proof}

% ================================================================
\section{Conceptual Mapping}\label{sec:mapping}

\begin{table*}[t]
\centering
\caption{Complete mapping between path tracing and agentic workflow concepts.}
\label{tab:mapping}
\small
\begin{tabular}{@{}llll@{}}
\toprule
\textbf{Path Tracing} & \textbf{Math} & \textbf{Agent Workflow} & \textbf{Math} \\
\midrule
Ray & $r(t) = \mathbf{o} + t\mathbf{d}$ & Single invocation & $a \sim p(\cdot \mid s)$ \\
Path & $\bar{x} = (x_0,\ldots,x_K)$ & Execution chain & $\bar{p} = (s_0,a_0,\ldots,s_K)$ \\
Radiance & $L(\bx, \omega)$ & Output quality & $Q(s)$ \\
BRDF & $f_r(\bx, \omega_i, \omega_o)$ & Transfer function & $T(s, a)$ \\
Rendering eq.\ & Eq.~\eqref{eq:rendering} & Quality eq.\ & Eq.~\eqref{eq:agent-quality} \\
SPP & $N$ rays/pixel & Samples/task & $N$ calls/task \\
Point light & $\mu = 0$ & Exact-answer task & $\mu(\Astar) = 0$ \\
NEE & Light sampling & Ref-guided gen.\ & RAG + tools \\
BSDF sampling & $\omega \sim f_r$ & Free generation & $a \sim \pLLM$ \\
MIS & Weighted combo & Multi-strategy & Weighted combo \\
Russian Roulette & Stochastic term.\ & Validator pruning & Quality-based stop \\
Importance Sampling & Guide toward $f$ & Constraint Buffer & Structured prompting \\
MLT & MCMC mutation & Doc refinement & Section-wise revision \\
BVH & Spatial index & Knowledge index & Vector DB \\
Environment map & Precomputed illum.\ & Reference KB & Curated exemplars \\
Denoising & Post-process & Result refinement & Synthesis agent \\
Energy conservation & $\int f_r \leq 1$ & Quality bound & $T(s,a) \leq 1$ \\
\bottomrule
\end{tabular}
\end{table*}

% ================================================================
\section{Experimental Design}\label{sec:experiments}

\subsection{Experiment 0: Point Light Validation}

\textbf{Hypothesis}: For near-zero-measure tasks,
NEE is necessary; pure sampling fails $\forall N$.

\begin{itemize}[nosep]
  \item \textbf{Task}: Strict JSON schema, precise API calls
  \item \textbf{Methods}: Pure sampling ($N \in \{1,10,100,1000\}$),
    NEE only, MIS
  \item \textbf{Models}: GPT-4o-mini, Claude Haiku, Gemini Flash
  \item \textbf{Metrics}: Exact match, schema pass rate
\end{itemize}

\subsection{Experiment 1: Basic Monte Carlo}

Validate $O(1/N)$ convergence on simplified Art Bible sections.
$N \in \{1,4,16,64,256\}$ across small vs.\ large models.
Selection via best-of-$N$ and majority voting.

\subsection{Experiment 2: Russian Roulette}

5-step Art Bible chain with RR at $Q_{\text{threshold}} \in \{0.3, 0.5, 0.7\}$.
Measure token consumption vs.\ quality.

\subsection{Experiment 3: Importance Sampling}

Compare unconstrained vs.\ Constraint Buffer generation.
Measure variance of quality scores and convergence speed.

\subsection{Experiment 4: MLT Mutation}

Compare full regeneration ($N\times$) vs.\
MLT ($1 + (N{-}1)$ mutations) on full Art Bible.
Metric: quality after $N$ total calls.

\subsection{Experiment 5: MIS Multi-Strategy}

Mixed-difficulty Art Bible.
Free only vs.\ RAG only vs.\ Tool only vs.\ MIS ($N/3$ each).

\subsection{Experiment 6: Sampling Efficiency}

Fixed budgets $B \in \{1, 5, 10, 50\}$ USD.
Small models (mini/Haiku/Flash) vs.\ large (4o/Sonnet/Pro).
Measure quality and empirical $\eta_M$.

% ================================================================
\section{Discussion}\label{sec:discussion}

\subsection{Where the Analogy Holds}

Both~\eqref{eq:rendering} and~\eqref{eq:agent-quality} are
Fredholm equations of the second kind.
All variance reduction techniques depend only on integral structure,
not physical interpretation.

\subsection{Where the Analogy Breaks}

\begin{itemize}[nosep]
  \item \textbf{Energy conservation}: LLMs can amplify errors ($T > 1$).
  \item \textbf{BRDF reciprocity}: No agent analogue.
  \item \textbf{Deterministic geometry}: Agent state transitions are stochastic.
  \item \textbf{Continuous space}: Agent output is mixed continuous/discrete.
  \item \textbf{Known lights}: Correct answers are usually unknown a priori.
\end{itemize}

\subsection{The Verifier as Shadow Ray}

The Validator Agent maps to the \emph{shadow ray}:
both are binary functions gating a sample's contribution.
The quality scorer maps to throughput measurement.

\subsection{Subsuming Existing Approaches}

Best-of-$N$ = basic MC;
Self-Consistency = MC + majority vote;
Tree-of-Thought = MCTS;
RAG = NEE;
Tool use = point light sampling;
Fork-Merge = parallel paths + MIS;
Sherlock = RR with counterfactual $q_k$.

The unifying perspective: \textbf{all are specific instances of
Monte Carlo variance reduction applied to the agent quality integral.}

% ================================================================
\section{Conclusion}\label{sec:conclusion}

We have established a rigorous mathematical mapping between
Monte Carlo light transport and LLM agentic workflows.
The Agent Quality Equation~\eqref{eq:agent-quality} provides
a unifying framework under which existing techniques
(Best-of-$N$, self-consistency, RAG, tool use)
emerge as special cases of well-studied variance reduction methods.

The Point Light Theorem (\cref{thm:unreachable,thm:nee-reach})
provides the sharpest result:
for tasks where correct answers occupy a zero-measure set,
no amount of random sampling suffices---NEE-style approaches
are a mathematical necessity.

We believe the rendering community's decades of experience
with stochastic methods represent an untapped source of algorithms
for the AI agent era.

% ================================================================
\bibliography{references}

\end{document}
